\documentclass[11pt, oneside]{article} 
\usepackage{geometry}
\geometry{letterpaper} 
\usepackage{graphicx}
	
\usepackage{amssymb}
\usepackage{amsmath}
\usepackage{parskip}
\usepackage{color}
\usepackage{hyperref}

\graphicspath{{/Users/telliott/Github-Math/figures/}}

\title{Similar triangles}
\date{}

\begin{document}
\maketitle
\Large

%[my-super-duper-separator]

\label{sec:similarity_and_ratios}

We've defined similar triangles as not congruent, but with all three angles equal.  Previously we showed that if we superimpose angles at one vertex, so the adjacent sides coincide, then the sides opposite that angle are parallel.

These ideas are closely related and easily proven in both directions.  Equal angles $\iff$ parallel sides.  We call this property of all angles equal and the third side parallel, similarity.

A third property which is associated with similar triangles is that the sides are all in the same proportion.

We now explore a third statement:  that similarity means having corresponding sides in the same ratio and vice-versa.  By corresponding, we mean pairs of sides that lie opposite to the same angle in the two triangles.

$\bullet$ \ $\parallel$ sides and equal $\angle \iff$ equal ratios

We will call these the \emph{forward} ratio theorem, and the \emph{converse} ratio theorem.

Note:  we included a proof of (the forward) theorem in the chapter on right triangles.  But this idea is so important in geometry that we explore it further here, including Euclid's proofs.

\subsection*{two sides similar}

In general, it is not enough to have two sides in the same ratio, the triangles may still not be similar.  Counterexamples are trivial to construct:  just change the angle between the two sides that are in the same ratio, for one of the triangles.

But if we have two sides in the same ratio, and we also know that the angle between those two sides is equal, then we have what you can think of as SAS for similarity.  This is a kind of hybrid of the two conditions, we know one pair of angles equal and two sides in proportion.  We look at this later in the chapter.
 
\subsection*{parallel third side implies equal ratios}

As background, we recall two fundamental ideas about triangles.  

The first is that if two triangles have their bases on the same line, and they also share the same vertex opposite, then they have exactly the same altitude.  It follows that the areas are in the same proportion as the lengths of the bases.  This is the \hyperref[sec:area_ratio_theorem]{\textbf{area-ratio theorem}}.  
\begin{center} \includegraphics [scale=0.5] {area11.png} \end{center}

There is only one altitude that can be drawn from a point to a base.  Since the triangles share the vertex point, they have the same altitude.

The second idea is that if two triangles share the same base, and the two opposing vertices both lie along a lie that is parallel to the base, then the triangles have the same area.  (See \hyperref[sec:triangle_area]{\textbf{here}}).
\begin{center} \includegraphics [scale=0.4] {area2.png} \end{center}
$\triangle PDE$ and $\triangle QDE$ have equal area and so does $\triangle PAB$.  \emph{Any} point along $PQR$ forms a triangle with base $DE$ that has the same area.  And any line segment with length equal to $DE$ anywhere along $A..J$ with a vertex on the line $PQR$ will have exactly the same area as well.

These two ideas can be combined.  If two triangles have their bases lying on the same line, and also the opposing vertices along a line parallel to the first, then their areas are in the same proportion as the lengths of the bases. 

\subsection*{colored triangles}

\label{sec:similarity_equal_ratios}

We first show a simplified proof of the forward theorem, with very few labels to avoid clutter.  The top panels below show slightly different views of the same two similar triangles, one shaded gray and a larger one outlined in red.  

The first triangle has sides $a$ and $b$ and the second has sides $A$ and $B$.  One of the third sides is black and the other red, and these lie parallel to each other.

In the left (top) panel we draw another triangle in magenta.  Now consider the combination of the magenta and gray together as a single triangle.  

\begin{center} \includegraphics [scale=0.4] {Euclid_VI_3d.png} \end{center}
If we angle our heads and consider $a$ as the \emph{base} of the gray triangle and its extension $A$ as the base of the combination, then the areas of the two triangles are in the same ratio as the bases, namely, $a/A$, as indicated in the lower part.  This follows from the area-ratio theorem, since they share the same point as the opposing vertex.

The same argument used for the cyan triangle yields the result given as $b/B$.

But the magenta and cyan triangles are equal in area.  The reason is that, viewed another way, they have the same base, the black horizontal line, and their vertices lie on the same line, namel y the red line running across the bottom, which is parallel to the bases of the magenta and cyan triangles.

Therefore the two ratios given above are equal:
\[ \frac{a}{A} = \frac{b}{B} \]

Equal angles (parallel third side) implies equal ratios.

$\square$

This result applies to any vertex of the two similar triangles and its adjacent sides, hence it applies to all three sides of the two triangles.

\subsection*{Euclid $VI-2$}

We repeat the same proof, with more labels.

\label{sec:Euclid6_2}

\begin{center} \includegraphics [scale=0.6] {Euclid_VI_2.png} \end{center}

In $\triangle ABC$, let $DE$ be drawn parallel to $BC$, with $AD$ in some ratio to $DB$, call it $k$.

Then, we claim that $AE$ is to $EC$ in the same ratio.

\emph{Proof.}

First, $\triangle BDE$ and $\triangle ADE$ have the same vertex $E$ and their bases, $AD$ and $DB$, lie on the same line, so they have the same altitude.  

\begin{center} \includegraphics [scale=0.4] {Euclid_VI_3a.png} \end{center}

Therefore, by the area-ratio theorem, their areas are in proportion to the lengths of the bases.

If we designate area as $\Delta$:
\[ \frac{\Delta_{ADE}}{\Delta_{BDE}} = \frac{AD}{BD} \]

Similarly, $\triangle CDE$ and $\triangle ADE$ have the same vertex $D$ and their bases, $AE$ and $CE$, lie on the same line (below).  Therefore they have the same altitude and their areas are in proportion to the lengths of the bases.  

\begin{center} \includegraphics [scale=0.4] {Euclid_VI_3b.png} \end{center}
\[ \frac{\Delta_{ADE}}{\Delta_{CDE}} = \frac{AE}{CE} \]

Combining these two results we have that
\[ \frac{\Delta_{CDE}}{\Delta_{BDE}} = \frac{AD}{BD} \cdot \frac{CE}{AE}  \]

Now we use the information about $DE \parallel BC$.

Triangles $\triangle BDE$ and $\triangle CDE$ have the same base, $DE$.

Because the points $B$ and $C$ lie on a line parallel to the base, their altitudes to that base have the same length, and therefore the triangles have the same area.

\begin{center} \includegraphics [scale=0.4] {Euclid_VI_3c.png} \end{center}

This means that in the previous expression, the left-hand side is equal to $1$, leaving
\[ \frac{AE}{CE} = \frac{AD}{BD} \]

$\square$

\subsection*{converse}

\label{sec:Euclid6_2_converse}

Euclid also runs this argument backward to prove the converse:  given equal ratios, it follows that $DE$ is parallel to $BC$.  

\emph{Proof}.

Given

\[ \frac{AD}{AB} = \frac{AE}{AC} = r \]

The ratio of areas of two triangles with the same altitude is in proportion to the bases.  Let that ratio be $r$.

\[ \frac{\Delta_{ADE}}{\Delta_{ABE}} = k = \frac{\Delta_{ADE}}{\Delta_{ACD}} \]

It follows that 
\[ \Delta_{ABE} = \Delta_{ACD} \]

and since they have a shared region $\triangle ADE$, the remainder of each one has the same property:
\[ \Delta_{BDE} = \Delta_{CDE} \]

Triangles with the same area and the same base have their top vertices lying along a line parallel to the base, so that they have the same height.

\[ DE \parallel BC \]

$\square$

Strictly speaking, we never provided a rigorous proof of the last statement which is the converse of what we said previously, but just gave it as a fact when we talked about triangular area.

\emph{Proof}.  (sketch)

Two triangles share the same base and have the same area but have different vertices.  The lines forming the altitudes for each triangle are perpendicular to the same line, so parallel to each other and also have the same length.  Therefore they form a rectangle, so the line holding the vertices is parallel to the base.

\subsection*{SAS for similar triangles}

\label{sec:SAS_similar}

Above, we showed that equal angles implies all sides with equal ratios, as well as the converse.  

Now, we mix and match the conditions, taking two sides in proportion and one angle shared.  This is often called SAS similarity.

We will prove that this mix of conditions are enough for two triangles to be similar, and then use that to give a more detailed proof of the converse theorem.

\begin{center} \includegraphics [scale=0.5] {mod_midpoint.png} \end{center}

\emph{Proof}.

Given two triangles with the same vertex angle at $A$ and two of three sides known to be in the same proportion ($BD/AB = CE/AC = k$).

Draw $EF$ parallel to $ABD$ and extend $BC$ to meet $EF$ at $F$.  

$\triangle ABC \sim \triangle CEF$ by vertical angles plus alternate interior angles (red dots).  The corresponding sides opposite the vertical angles are $EF$ and $AB$.

By the forward theorem
\[ \frac{EF}{AB} = k = \frac{CE}{AC} \]
but
\[ \frac{CE}{AC}  = \frac{BD}{AB} \]

Therefore, $BD = EF$.  We are given $BD \parallel EF$.  Therefore, $BDEF$ has one pair of opposing sides equal and parallel, so it is a parallelogram.  

It follows that $BC \parallel DE$ and $BF = DE$.

It also follows that $\angle D = \angle ABC$ by alternate interior angles, so we have two angles equal which means that $\triangle ABC \sim \triangle ADE$.  

The forward theorem then gives $(DE-BC)/BC = BD/AB = k$

$\square$

We can also use the SAS similarity theorem to prove the converse theorem, that equal ratios implies equal angles.  Consider these two nested triangles.
\begin{center} \includegraphics [scale=0.4] {converse_sim_2.png} \end{center}

We are given equal ratios (left panel), with
\[ \frac{a}{a'} = \frac{b}{b'} = \frac{c}{c'} \]
$c'$ is the difference between the whole length of the bottom and $c$.

Now draw a small triangle by extending side $c$ with $y$ such that $y = c'$.  Then $c$ is in the same ratio to $y$ as $b$ is to $b'$ and there is a shared vertical angle.  

Therefore, the new triangle is similar to the one with sides $a,b$ and $c$ by SAS similarity.  That means
\[ \frac{a}{x} = \frac{b}{b'} = \frac{c}{y} \]

Comparing that to what we were given, we have that $x$ is equal to the side opposite, namely, $a'$.  In the quadrilateral, we have both pairs of opposing sides equal.  So the additional triangle forms part of a parallelogram.

Similar triangles also gives $x \parallel a+a'$ by alternate interior angles.  Thus, we also have one pair of opposing sides parallel.  

Since $x \parallel a'$, the angle at the lower left is equal to the opposing vertex between $x$ and $y$, which is the third angle we needed to prove equal.  But we already knew that, by the triangle sum theorem.

$\square$

Two other proofs relating similarity and equal ratios are in the appendix.  One relies on a subtle idea from calculus called a limit:  \hyperref[sec:similarity_theorem]{\textbf{AAA similarity theorem}} (Kiselev).  

We also showed that this follows for right triangles easily from the properties of such triangles (\hyperref[sec:similar_right_triangles]{\textbf{here}}), and is related to the area-ratio theorem.

Two additional proofs are \hyperref[sec:SAS_similar_cosine]{\textbf{here}}.

\subsection*{important note}
To emphasize the importance of the reverse theorem, we include here a re-scaling proof of the Pythagorean theorem.

\subsection*{scaled triangle proof}

\label{sec:Pythagoras_scaled_triangles}

\emph{Proof.}

Draw a right triangle and label sides $a,b$ and $c$.

\begin{center} \includegraphics [scale=0.5] {pyth1.png} \end{center}
Now, flip and rotate the triangle.  Make a copy of that one and rotate the copy so that the shortest side is to the left.  Enlarge the copy until the two adjacent sides are equal in length.
\begin{center} \includegraphics [scale=0.5] {pyth2.png} \end{center}

These are similar triangles so the angles are all equal, but the sides are proportional, multiplied by a common factor, which here is $b$.

The adjacent sides show that $ab = b$, so this choice of scaling defines $a = 1$.  And since $a = 1$, we can also view each side of the original triangle as being scaled by a factor of $a$.  

We just tried scaling by a factor of $b$ and then by $a$, so that naturally suggests we try scaling by a factor of $c$.  

\begin{center} \includegraphics [scale=0.5] {pyth2b.png} \end{center}
All sides of the enlarged triangle have been multiplied by the same factor, namely, $c$.  Notice that $c = ac$, as required.

We're nearly done.  Put the three triangles together.

\begin{center} \includegraphics [scale=0.5] {pyth2c.png} \end{center}

The lower left and right corners are places where two vertices come together.  By complementary angles, these sum to be right angles.  

Since there are two other right angles at the vertices of the quadrilateral, it is a rectangle.  Therefore the place where three vertices come together is a straight line, which we can also verify because there are two complementary angles plus a right angle.

Since the composite figure is a rectangle, opposing sides are equal, so $a^2 + b^2 = c^2$.

$\square$

\subsection*{problem}

Draw the line connecting the vertices and draw diagonals to show that the resulting quadrilateral is a rectangle.

\subsection*{problem}

\begin{center} \includegraphics [scale=0.5] {similar23.png} \end{center}

The problem states that the bases are parallel ($XY \parallel BC$) and also that the subdivision produces \emph{equal areas}.  We are asked to find the ratio $AX:XB$.

\emph{Solution}.

Recall from a previous chapter that for two similar triangles, the altitudes to corresponding sides are in the same ratio as any of the three pairs of sides themselves.  

The altitudes $h$ and $H$ to $BC$ (not drawn) are also in the same ratio as the sides, so let $H = kh$ --- $H$ lies on $BC$.

Then twice the area of the top triangle is $AX \cdot h$ and twice the area of the whole is $AB \cdot H = k AB \cdot h$

We have that the whole is twice the smaller area so the ratio is equal to $2$:
\[ \frac{k AB \cdot h}{AX \cdot h} = 2 \]
But $AB/AX$ is also equal to $k$ so we have that $k^2 = 2$ and $k = \sqrt{2}$.

However, we are not asked for $k$.  Instead, we want the ratio of $AX$ to the smaller piece along the bottom. Let $AX = a$.  Then the whole is $A = ka$.  The difference is the small piece, $XB = A - a = a(k - 1)$.  We need
\[ \frac{a}{a(k - 1)} = \frac{1}{k - 1}  = \frac{1}{\sqrt{2} - 1} \]

\subsection*{problem}

The figure below is from Acheson's wonderful book aptly titled \emph{The Wonder Book of Geometry}.  He shows this problem, which he says "[goes] back to at least AD 850, when it appeared in a textbook by the Indian mathematician Mahavira."

\begin{center} \includegraphics [scale=0.5] {Acheson_ladders.png} \end{center}

Looking down an alleyway, you see two ladders arranged as shown and wonder about the point where they cross, at height $h$ and distances $c_1$ and $c_2$ from the edges of the alley, where the width of the alley is $c = c_1 + c_2$.

By similar triangles
\[ \frac{c_1}{h} = \frac{c}{b} \]

Can you see why?

Going the opposite direction
\[ \frac{c_2}{h} = \frac{c}{a} \]

Adding the two equations and substituting for $c_1 + c_2$:
\[ \frac{c}{h} = \frac{c}{a} + \frac{c}{b} \]

Thus
\[ \frac{1}{h} = \frac{1}{a} + \frac{1}{b} \]

That's a simple and interesting result.  $h$ depends only on $a$ and $b$ and not on $c, c_1$, or $c_2$.

If you think about it, you should see that in the previous problem we never used the information that the sides and the height were vertical, only parallel.

\begin{center} \includegraphics [scale=0.4] {similar25.png} \end{center}
Work through the same proof for the figure above to show that $1/x + 1/y = 1/z$.

\subsection*{problem}

Here's a problem from the web.  Given that $AC \parallel EF$ and $AB \parallel DF$.

We are to prove that the sum of the altitudes of the small triangles is equal to the altitude of the large one.

\begin{center} \includegraphics [scale=0.4] {prob_similar_tri2.png} \end{center}

Informal solution:  The smaller triangles are all similar to $\triangle ABC$ by the alternate interior angles and vertical angle theorems.

For similar triangles, not only are the sides in the same ratio to each other, but so are other measures like the altitudes to a particular side.  So if we label the bases $b_1$ etc., collectively $b_i$ , then we have that
\[ \frac{b_i}{h_i} = \frac{b}{h} \]
\[ b_i = b \cdot h_i/h \]
for each of the $b_i$.

But the sum of the $b_i$ is simply equal to $b$ so
\[ b_1 + b_2 + b_3 = b \cdot (h_1/h + h_2/h + h_3/h) \]
\[ b = b \cdot (h_1/h + h_2/h + h_3/h) \]
\[ 1 = h_1/h + h_2/h + h_3/h \]
\[ h = h_1 + h_2 + h_3 \]

\subsection*{problem}

Given two triangles that are not similar, where the ratio $AD/AB = r$ and $AE/AC = s$.  We are asked to show that 

\[ \frac{\triangle_{ADE}}{\triangle_{ABC}} = rs \] 

\begin{center} \includegraphics [scale=0.4] {similarity_by_area2.png} \end{center}

\emph{Solution}.

Draw $DF \parallel BC$.  Now $\triangle ADF \sim \triangle ABC$ so $AF/AC = r$.

By our previous result
\[ \frac{\triangle_{ADF}}{\triangle_{ABC}} = r^2 \] 

Since $AE/AC = s$ and $AF/AC = r$, $AE/AF = s/r$.  As triangles with a common vertex and bases in that proportion, these areas are in the same proportion:

\[ \frac{\triangle_{ADE}}{\triangle_{ADF}} = \frac{s}{r} \] 

Multiply the two ratios together to obtain

\[ \frac{\triangle_{ADE}}{\triangle_{ABC}} = sr \] 

$\square$

\emph{Solution}.  (alternate).

\begin{center} \includegraphics [scale=0.4] {similarity_by_area2.png} \end{center}

Looking ahead to trigonometry, in any triangle, twice the area can be computed as the product of two sides flanking an angle times the \emph{sine} of the angle.  (The reason is that the altitude to one side, divided by the length of the second side, is defined to be the sine of the angle.)

\[ 2 (ADE) = AD \cdot AE \cdot \sin \angle DAE \]
\[ 2 (ABC) = AB \cdot AC \cdot \sin \angle BAC \]

But $\angle DAE = \angle BAC$, so they have the same sine, and the ratio of areas is just
\[ \frac{AD \cdot AE}{AB \cdot AC} = rs \]

$\square$

To rework this proof in terms of familiar concepts, draw the altitude from $D$ to $AE$, and also the one from $B$ to $AC$

They form similar right triangles including the angle at $A$.  The altitudes scale like $AD/AB = r$.  But the bases scale like $AE/AC = s$.  

And area scales like the product of the two, namely, $rs$.

\subsection*{Centroid}
Consider the $\triangle ABC$. Draw the medians $BE$ and $CF$ (bisectors of the sides).  Extend $AG$ through the intersection of the two medians at $G$ to $D$ and then finally to $H$, where $CH$ is drawn parallel to $BGE$.  We profess not to know anything about $BH$ yet.

\begin{center} \includegraphics [scale=0.3] {centroid_pgram1.png} \end{center}

Since $CH \parallel BGE$ and $AE = EC$, we have that $\triangle AGE \sim \triangle AHC$ with ratio $2$ by the midpoint theorem (or what we've called SAS similarity).  Thus, $AG = GH$.

But we also have $AF = FB$, which means that $\triangle AFG \sim \triangle ABH$ with ratio $2$.  This gives $BH \parallel FGC$.

Therefore, $BGCH$ has two pairs of opposing sides parallel so it is a parallelogram.  The diagonals cross at their midpoints, which means that $AD$ is a median ($BC$ is bisected at $D$).  

In addition, $GD = DH$.  Thus, $GD$ is one-half of $GH$, and $GH$ is equal to $AG$, so $GD$ is one-quarter of $AH$ and one-third of $AD$.  $G$ lies on all three medians, and is called the \emph{centroid} of the triangle.  For a physical triangle it would be the center of mass.

\subsection*{pyramid height}

As we said earlier, Thales was from Miletus and he lived around 600 BC.  Thales is believed to have traveled extensively and was likely of Phoenician heritage.  As you probably know, the Phoenicians were famous sailors who founded many settlements around the Mediterranean.  

They competed with the mainland Greeks and later with the Romans for colonies, and their major city, Carthage, was destroyed much later by the Romans, in the third Punic War.  Hannibal rode his famous elephants over the Alps in the second Punic war.

During his travels, Thales went to Egypt, home to the great pyramids at Giza, which were already ancient then.  They had been built about 2560 BC (dated by reference to Egyptian kings) and were already 2000 years old at that time!

The story is that Thales asked the Egyptian priests about the height of the Great Pyramid of Cheops, and they would not tell him.  So he set about measuring it himself.  He used similar triangles.  I'm sure he wrote down his answer, but I'm not aware that it survives.  The current height is 480 feet.

\begin{center} \includegraphics [scale=0.25] {Thales_theorem_6.png} \end{center}

\end{document}